% 主旨：三组同代码同超参的冷启动对比中，用掩码的配置能推塔、不用掩码的不能；混合奖励在掩码之上
%       提高终点指标。续训模型的六局对战与离线固定权重实验作补充；最后分析失败模式。
\section{实验}
\label{sec:experiments}

\subsection{实验设置}
\label{sec:setup}

云端实验全部在开悟 V63.0.9 的分布式训练集群上进行，训练与评估协议见 \cref{sec:protocol}。\cref{tab:configs} 列出六组配置。

“仅奖励”“仅掩码”“完整”三组使用同一份代码和 \cref{sec:base} 的超参数，只切换掩码开关或奖励权重，均从随机初始化训练 \qty{30}{\minute}，是本文的受控对比。“基线”为平台官方配置，网络与 PPO 损失与本文相同，但没有掩码和冷启动偏置，学习率在 \num{5000} 步内由 \num{1e-3} 降至 \num{1e-4}，熵系数为 \num{0.025} 且不按分支缩放，训练 \qty{25}{\minute}；它只作参照，不参与受控对比。“完整（旧规则）”采用同一方法的较早规则与权重：规则中尚无受压选敌与回泉保持，走廊血线为 \num{0.40}，英雄血量与补刀权重分别为 \num{0.80} 与 \num{0.30}。它与“完整”不是同一配置，不能视为重复实验。“续训”从“完整（旧规则）”的终点检查点出发，换用当前规则与权重再训练 \qty{60}{\minute}。各组对应的平台任务编号见\cref{sec:appendix}。

\begin{table}[htbp]
  \centering
  \caption{云端实验配置}
  \label{tab:configs}
  \begin{tabular}{lccccr}
    \toprule
    配置 & 战术掩码 & 奖励 & 冷启动偏置 & 初始化 & 时长 / min \\
    \midrule
    基线 & 无 & 稀疏 & 无 & 随机 & 25 \\
    仅奖励 & 无 & 混合 & 有 & 随机 & 30 \\
    仅掩码 & 有 & 稀疏 & 有 & 随机 & 30 \\
    完整 & 有 & 混合 & 有 & 随机 & 30 \\
    完整（旧规则） & 旧版 & 旧权重 & 有 & 随机 & 30 \\
    续训 & 有 & 混合 & -- & 完整（旧规则）终点 & 60 \\
    \bottomrule
  \end{tabular}
\end{table}

指标取训练结束时平台监控面板的最后一个统计点；监控按评估局滚动汇总，按 \cref{sec:protocol} 的协议，评估局约占全部对局的十一分之一。每组配置只运行一次，没有不同随机种子的重复，这是本文结果最主要的不确定来源（\cref{sec:limitations}）。

\subsection{冷启动对比}
\label{sec:ablation}

\cref{tab:ablation} 给出各组训练终点的监控指标，\cref{fig:ablation} 给出训练过程。用了掩码的三组都能推塔，没用掩码的两组一局未胜。“完整”在第 \qty{1}{\minute} 的统计点上胜率为 \num{0.71}，第 \qty{12}{\minute} 起保持在 \numrange{0.96}{1.00}，终点敌方外塔只剩 \num{38} 点生命值。“完整（旧规则）”的终点胜率为 \num{0.96}，另一套规则与权重下同样能推塔。

\begin{table}[htbp]
  \centering
  \caption{冷启动训练终点对 \texttt{common\_ai} 的监控指标（外塔满值 \num{12000}）}
  \label{tab:ablation}
  \begin{tabular*}{\linewidth}{@{\extracolsep{\fill}}lrrrcrr@{}}
    \toprule
    配置 & 胜率 & 敌方外塔余血 & 己方外塔余血 & K/D & 伤抗比 & 局长 / 帧 \\
    \midrule
    基线 & 0.00 & 12000 & 52 & 0.00/0.25 & 0.15 & 7381 \\
    仅奖励 & 0.00 & 12000 & 4 & 0.00/0.00 & 0.00 & 7169 \\
    仅掩码 & 0.80 & 843 & 3526 & 2.53/1.73 & 0.95 & 11790 \\
    完整 & 1.00 & 38 & 8107 & 3.76/1.16 & 1.30 & 11112 \\
    完整（旧规则） & 0.96 & 274 & 7874 & 2.58/1.17 & 0.95 & 9603 \\
    \bottomrule
  \end{tabular*}
\end{table}

\begin{figure}[htbp]
  \centering
  \includegraphics[width=\linewidth]{cloud_ablation.pdf}
  \caption{冷启动训练过程中对 \texttt{common\_ai} 的监控指标。横轴以首个训练步数监控点为零点（其前的胜率统计点落在负半轴）；“基线”只训练 \qty{25}{\minute}，其曲线在各图中与“仅奖励”几乎重合。}
  \label{fig:ablation}
\end{figure}

去掉掩码、保留混合奖励与冷启动偏置的“仅奖励”与官方基线一样一局未胜，敌方外塔始终满血，己方外塔几乎被推平。它在评估局中对敌方英雄的伤害、击杀和死亡都为零，没有形成有效的对英雄输出。

反过来，保留掩码、只给稀疏推塔奖励的“仅掩码”在第一个统计点就有 \num{0.67} 的胜率，此后在 \numrange{0.67}{0.83} 之间波动，终点为 \num{0.80}。这部分胜率从训练开始就存在，应主要来自掩码规则本身：推进走廊把英雄送向敌塔，射程门控让刚初始化的策略也只对射程内的目标出手。与“完整”相比，它的终点己方外塔余血为 \num{3526}（“完整”为 \num{8107}），场均死亡 \num{1.73}（\num{1.16}），伤抗比 \num{0.95}（\num{1.30}）。

保留掩码时，混合奖励把终点胜率由 \num{0.80} 提到 \num{1.00}，己方外塔余血由 \num{3526} 提到 \num{8107}，场均死亡由 \num{1.73} 降到 \num{1.16}；去掉掩码时，同样的混合奖励没有带来胜局。掩码决定了策略能否推塔，混合奖励则在掩码之上改善了生存与防守。

\subsection{续训与模型对战}
\label{sec:h2h}

续训 \qty{60}{\minute} 后，监控中对 \path{common_ai} 的胜率在全部 \num{64} 个统计点上保持 \num{1.00}，终点敌方外塔余血 \num{54.48}、己方 \num{9895.19}，K/D 为 \num{4.35}/\num{1.00}，伤抗比 \num{1.577}；同一指标在续训起点“完整（旧规则）”上分别为 \num{273.88}、\num{7873.50}、\num{2.58}/\num{1.17} 与 \num{0.945}。

为在规则对手之外比较两个冻结模型，用平台的模型对战功能让续训终点（记为 A）与续训起点（记为 B）进行三轮换边、共六局确定性对战。结果见 \cref{tab:h2h} 与 \cref{fig:h2h}：A 胜 \num{4} 局、B 胜 \num{2} 局，无超时与异常，A 在红蓝两侧均为 \num{2} 胜 \num{1} 负。A 的场均 KDA 为 \num{1.22}，B 为 \num{0.58}；A 的场均经济高出 \num{623}。六局中 A 对英雄的伤害均高于 B。

\begin{table}[htbp]
  \centering
  \caption{续训终点（A）与续训起点（B）的六局换边对战}
  \label{tab:h2h}
  \begin{tabular*}{\linewidth}{@{\extracolsep{\fill}}cccrrrr@{}}
    \toprule
    局 & A 阵营 & 胜方 & 局长 / 帧 & A 的 K/D & A 经济 & B 经济 \\
    \midrule
    1 & 红 & A & 9488 & 0/1 & 4862 & 4352 \\
    2 & 红 & B & 15038 & 2/3 & 8503 & 8962 \\
    3 & 蓝 & A & 8858 & 2/0 & 4702 & 3292 \\
    4 & 蓝 & A & 10088 & 2/1 & 5771 & 4440 \\
    5 & 蓝 & B & 8036 & 1/2 & 3240 & 3997 \\
    6 & 红 & A & 10796 & 4/0 & 6517 & 4814 \\
    \midrule
    平均 & -- & 4:2 & 10384 & 1.83/1.17 & 5599 & 4976 \\
    \bottomrule
  \end{tabular*}
\end{table}

\begin{figure}[htbp]
  \centering
  \includegraphics[width=\linewidth]{head_to_head.pdf}
  \caption{六局换边对战。(a) 逐局英雄伤害，柱上标注胜方与双方 K/D；(b) 逐局经济；(c) 六局场均指标中 A 相对 B 的倍数，括号内为 A/B 原值，死亡取 B/A。}
  \label{fig:h2h}
\end{figure}

若各局独立且两模型每局胜率均为 \num{0.5}，A 至少赢 \num{4} 局的单侧二项概率为 \num{0.34}，因此六局结果还不能区分两者的实力；A 在六局中的伤害与平均经济均更高。

\subsection{离线固定权重实验}
\label{sec:local}

云端对比每次都重新训练，掩码的作用与参数更新交织在一起。为单独观察掩码，在云端实验之前用离线引擎做了两组固定权重的探索性对比。离线环境是平台提供的独立 GameCore 引擎，观测为 \num{705} 维扁平向量，策略为两层 \num{256} 维全连接网络，判胜规则与对手（\path{common_ai}）同云端一致；掩码是 \cref{tab:mask-rules} 中补给与走廊两组规则的早期版本。结果见 \cref{tab:local}，括号中的 $p$ 值均为双侧 Fisher 精确检验，未做多重比较校正。

\begin{table}[htbp]
  \centering
  \caption{离线引擎上的固定权重对比}
  \label{tab:local}
  \begin{tabular}{llrrr}
    \toprule
    权重 & 决策方式或掩码 & 胜 / 局 & 场均死亡 & 英雄伤害 \\
    \midrule
    权重 1 & 随机采样 & 5/20 & 2.25 & 20058 \\
    权重 1 & $\argmax$ & 0/20 & 1.90 & 6234 \\
    \midrule
    权重 2 & 无补给规则 & 1/10 & 3.30 & 13906 \\
    权重 2 & 加回泉 & 3/5 & 0.00 & 30166 \\
    权重 2 & 加回泉、血包与推进走廊 & 8/10 & 0.40 & 16118 \\
    \bottomrule
  \end{tabular}
\end{table}

第一组中，同一权重在随机采样下赢 \num{5}/\num{20}，取 $\argmax$ 后为 \num{0}/\num{20}（$p=0.047$），英雄伤害降为约三分之一，回放中 $\argmax$ 策略反复空放技能、在无目标处游荡。这说明对该检查点而言，随机采样掩盖了确定性策略的缺陷；而正式评估使用确定性动作，不可取的动作更适合在分布上直接移除，而不是依赖采样随机性绕开。

第二组固定网络权重，只改掩码。加入回泉后胜局由 \num{1}/\num{10} 变为 \num{3}/\num{5}，再加入血包与推进走廊后为 \num{8}/\num{10}（与 \num{1}/\num{10} 相比 $p=0.005$），场均死亡由 \num{3.30} 降至 \num{0.40}。网络没有更新，表现的变化来自掩码对动作的约束，这与云端“仅掩码”在训练开始时即有较高胜率一致。

\subsection{失败模式}
\label{sec:failure}

“仅奖励”的失败方式值得单独说明。它在自对弈中的平均局长为 \num{14789} 帧，比“完整”长约 \qty{40}{\percent}，每帧经济只有“完整”的 \qty{60}{\percent}；对 \texttt{common\_ai} 则既没有对英雄输出，也没有阵亡。一种解释是：混合奖励中的血量、阵亡和塔伤三项都是负向信号，未经约束的早期探索一旦接战，多半先承受这些惩罚。推塔与击杀的正向回报则需要连续多步的正确动作才能拿到，策略于是停在回避交战的局部最优。掩码删去了空动作和射程外的攻击，并在安全时把移动方向锁向敌塔，跳过了这段代价很高的探索。这一解释与上述指标一致，尚未经单独实验检验。

掩码中的回泉保持与受压选敌（\cref{tab:mask-rules}）来自对“完整（旧规则）”的回放统计。用其终点检查点重放三局（一胜两负、场均死亡 \num{2.00}），英雄场均 \num{17.7} 次转向泉水，其中 \num{10.0} 次半途折返；受到敌方英雄攻击时，仍有 \num{246} 步在攻击小兵或防御塔。两条规则分别针对这两种失败模式。

续训模型在对战中的两场负局暴露了剩余的问题。第 \num{2} 局持续 \num{15038} 帧，A 对英雄的伤害为 \num{43734}，高于 B 的 \num{36578}，但阵亡 \num{3} 次；第 \num{5} 局在 \num{8036} 帧结束，A 的经济以 \num{3240} 比 \num{3997} 落后。对 \path{common_ai} 的场均死亡也仍为 \num{1.00}。从这两局看，问题在于生存与推进节奏而不在输出，具体是撤退血线还是推塔时机，需要逐帧轨迹分析确定。
